% Diagrama de casos de uso — incluído em textos.tex por % Diagrama de casos de uso — incluído em textos.tex por % Diagrama de casos de uso — incluído em textos.tex por % Diagrama de casos de uso — incluído em textos.tex por \input{figuras/casos-de-uso}
% Requer no preâmbulo de principal.tex: \usetikzlibrary{shapes,positioning}
\begin{figure}[htbp]
    \centering
    \caption{Diagrama de casos de uso da aplicação}
    \label{fig:casos-de-uso}
    \vspace{-0.2cm}
    \begin{tikzpicture}[
        caso/.style={
            draw=black,
            thick,
            ellipse,
            font=\small,
            inner sep=0pt,
            minimum height=0.9cm
        },
        assoc/.style={thick},
        ator/.pic={
            \draw[thick] (0,1.05) circle (0.2);
            \draw[thick] (0,0.85) -- (0,0.3);
            \draw[thick] (-0.35,0.65) -- (0.35,0.65);
            \draw[thick] (0,0.3) -- (-0.3,-0.15);
            \draw[thick] (0,0.3) -- (0.3,-0.15);
        }
    ]
        \draw[thick] (-3.9,4.0) rectangle (3.9,-4.9);
        \node[font=\small\bfseries, anchor=north] at (0,3.95) {Aplicação web};

        \node[caso] (uc1) at (0,2.9) {Consultar indicadores gerais};
        \node[caso] (uc2) at (0,1.7) {Filtrar estabelecimentos};
        \node[caso] (uc3) at (0,0.5) {Explorar o mapa de pontos};
        \node[caso] (uc4) at (0,-0.7) {Consultar densidade por área};
        \node[caso] (uc5) at (0,-1.9) {Buscar na tabela por nome};
        \node[caso] (uc6) at (0,-3.1) {Ver data e nível de coordenada};
        \node[caso] (uc7) at (0,-4.3) {Carregar nova fotografia};

        \pic at (-5.8,-0.6) {ator};
        \node[font=\small, below] at (-5.8,-0.75) {Visitante};
        \coordinate (v) at (-5.4,0.35);
        \foreach \uc in {uc1,uc2,uc3,uc4,uc5,uc6} {
            \draw[assoc] (v) -- (\uc.west);
        }

        \pic at (5.8,-4.9) {ator};
        \node[font=\small, below] at (5.8,-5.05) {Administrador};
        \draw[assoc] (5.4,-4.15) -- (uc7.east);
    \end{tikzpicture}
    \vspace{-0.2cm}
    \\\textbf{\footnotesize Fonte: Elaborado pelo autor.}
\end{figure}

% Requer no preâmbulo de principal.tex: \usetikzlibrary{shapes,positioning}
\begin{figure}[htbp]
    \centering
    \caption{Diagrama de casos de uso da aplicação}
    \label{fig:casos-de-uso}
    \vspace{-0.2cm}
    \begin{tikzpicture}[
        caso/.style={
            draw=black,
            thick,
            ellipse,
            font=\small,
            inner sep=0pt,
            minimum height=0.9cm
        },
        assoc/.style={thick},
        ator/.pic={
            \draw[thick] (0,1.05) circle (0.2);
            \draw[thick] (0,0.85) -- (0,0.3);
            \draw[thick] (-0.35,0.65) -- (0.35,0.65);
            \draw[thick] (0,0.3) -- (-0.3,-0.15);
            \draw[thick] (0,0.3) -- (0.3,-0.15);
        }
    ]
        \draw[thick] (-3.9,4.0) rectangle (3.9,-4.9);
        \node[font=\small\bfseries, anchor=north] at (0,3.95) {Aplicação web};

        \node[caso] (uc1) at (0,2.9) {Consultar indicadores gerais};
        \node[caso] (uc2) at (0,1.7) {Filtrar estabelecimentos};
        \node[caso] (uc3) at (0,0.5) {Explorar o mapa de pontos};
        \node[caso] (uc4) at (0,-0.7) {Consultar densidade por área};
        \node[caso] (uc5) at (0,-1.9) {Buscar na tabela por nome};
        \node[caso] (uc6) at (0,-3.1) {Ver data e nível de coordenada};
        \node[caso] (uc7) at (0,-4.3) {Carregar nova fotografia};

        \pic at (-5.8,-0.6) {ator};
        \node[font=\small, below] at (-5.8,-0.75) {Visitante};
        \coordinate (v) at (-5.4,0.35);
        \foreach \uc in {uc1,uc2,uc3,uc4,uc5,uc6} {
            \draw[assoc] (v) -- (\uc.west);
        }

        \pic at (5.8,-4.9) {ator};
        \node[font=\small, below] at (5.8,-5.05) {Administrador};
        \draw[assoc] (5.4,-4.15) -- (uc7.east);
    \end{tikzpicture}
    \vspace{-0.2cm}
    \\\textbf{\footnotesize Fonte: Elaborado pelo autor.}
\end{figure}

% Requer no preâmbulo de principal.tex: \usetikzlibrary{shapes,positioning}
\begin{figure}[htbp]
    \centering
    \caption{Diagrama de casos de uso da aplicação}
    \label{fig:casos-de-uso}
    \vspace{-0.2cm}
    \begin{tikzpicture}[
        caso/.style={
            draw=black,
            thick,
            ellipse,
            font=\small,
            inner sep=0pt,
            minimum height=0.9cm
        },
        assoc/.style={thick},
        ator/.pic={
            \draw[thick] (0,1.05) circle (0.2);
            \draw[thick] (0,0.85) -- (0,0.3);
            \draw[thick] (-0.35,0.65) -- (0.35,0.65);
            \draw[thick] (0,0.3) -- (-0.3,-0.15);
            \draw[thick] (0,0.3) -- (0.3,-0.15);
        }
    ]
        \draw[thick] (-3.9,4.0) rectangle (3.9,-4.9);
        \node[font=\small\bfseries, anchor=north] at (0,3.95) {Aplicação web};

        \node[caso] (uc1) at (0,2.9) {Consultar indicadores gerais};
        \node[caso] (uc2) at (0,1.7) {Filtrar estabelecimentos};
        \node[caso] (uc3) at (0,0.5) {Explorar o mapa de pontos};
        \node[caso] (uc4) at (0,-0.7) {Consultar densidade por área};
        \node[caso] (uc5) at (0,-1.9) {Buscar na tabela por nome};
        \node[caso] (uc6) at (0,-3.1) {Ver data e nível de coordenada};
        \node[caso] (uc7) at (0,-4.3) {Carregar nova fotografia};

        \pic at (-5.8,-0.6) {ator};
        \node[font=\small, below] at (-5.8,-0.75) {Visitante};
        \coordinate (v) at (-5.4,0.35);
        \foreach \uc in {uc1,uc2,uc3,uc4,uc5,uc6} {
            \draw[assoc] (v) -- (\uc.west);
        }

        \pic at (5.8,-4.9) {ator};
        \node[font=\small, below] at (5.8,-5.05) {Administrador};
        \draw[assoc] (5.4,-4.15) -- (uc7.east);
    \end{tikzpicture}
    \vspace{-0.2cm}
    \\\textbf{\footnotesize Fonte: Elaborado pelo autor.}
\end{figure}

% Requer no preâmbulo de principal.tex: \usetikzlibrary{shapes,positioning}
\begin{figure}[htbp]
    \centering
    \caption{Diagrama de casos de uso da aplicação}
    \label{fig:casos-de-uso}
    \vspace{-0.2cm}
    \begin{tikzpicture}[
        caso/.style={
            draw=black,
            thick,
            ellipse,
            font=\small,
            inner sep=0pt,
            minimum height=0.9cm
        },
        assoc/.style={thick},
        ator/.pic={
            \draw[thick] (0,1.05) circle (0.2);
            \draw[thick] (0,0.85) -- (0,0.3);
            \draw[thick] (-0.35,0.65) -- (0.35,0.65);
            \draw[thick] (0,0.3) -- (-0.3,-0.15);
            \draw[thick] (0,0.3) -- (0.3,-0.15);
        }
    ]
        \draw[thick] (-3.9,4.0) rectangle (3.9,-4.9);
        \node[font=\small\bfseries, anchor=north] at (0,3.95) {Aplicação web};

        \node[caso] (uc1) at (0,2.9) {Consultar indicadores gerais};
        \node[caso] (uc2) at (0,1.7) {Filtrar estabelecimentos};
        \node[caso] (uc3) at (0,0.5) {Explorar o mapa de pontos};
        \node[caso] (uc4) at (0,-0.7) {Consultar densidade por área};
        \node[caso] (uc5) at (0,-1.9) {Buscar na tabela por nome};
        \node[caso] (uc6) at (0,-3.1) {Ver data e nível de coordenada};
        \node[caso] (uc7) at (0,-4.3) {Carregar nova fotografia};

        \pic at (-5.8,-0.6) {ator};
        \node[font=\small, below] at (-5.8,-0.75) {Visitante};
        \coordinate (v) at (-5.4,0.35);
        \foreach \uc in {uc1,uc2,uc3,uc4,uc5,uc6} {
            \draw[assoc] (v) -- (\uc.west);
        }

        \pic at (5.8,-4.9) {ator};
        \node[font=\small, below] at (5.8,-5.05) {Administrador};
        \draw[assoc] (5.4,-4.15) -- (uc7.east);
    \end{tikzpicture}
    \vspace{-0.2cm}
    \\\textbf{\footnotesize Fonte: Elaborado pelo autor.}
\end{figure}
